\begingroup
\renewcommand{\arraystretch}{1.45}
\fancyhead[L]{\small BESTIMMUNG DER ANZIEHUNG}
\Needspace{44\baselineskip}
\AttractionArticle{16}

Was die Größen $P,Q$ betrifft, so wird offenbar jede von ihnen
\[
\begin{aligned}
&=\frac{1}{2(G+G')^{\frac32}}
\end{aligned}
\]
sooft $G'=G''$ gilt; in allen übrigen Fällen sind sie dagegen den transzendenten Größen zuzurechnen. Wie sich diese durch Reihen ausdrücken lassen, ist hinreichend bekannt. Wir hoffen jedoch, dass es den Lesern willkommen sein wird, wenn wir bei dieser Gelegenheit die Bestimmung dieser und anderer transzendenter Größen durch einen besonderen, höchst bequemen Algorithmus erläutern, den wir schon seit vielen Jahren häufig benutzt haben und den wir an anderer Stelle ausführlicher behandeln wollen.

\Needspace{36\baselineskip}
Es seien $m,n$ zwei positive Größen, und wir setzen
\[
\begin{aligned}
&m'=\frac12(m+n),\qquad n'=\surd{mn}
\end{aligned}
\]
sodass $m',n'$ das arithmetische beziehungsweise geometrische Mittel zwischen $m$ und $n$ sind. Wir werden voraussetzen, dass das geometrische Mittel stets positiv genommen wird. Ebenso sei
\[
\begin{aligned}
&m''=\frac12(m'+n'),\qquad n''=\surd{m'n'}\\[8pt]
&m'''=\frac12(m''+n''),\qquad n'''=\surd{m''n''}
\end{aligned}
\]
und so fort. Auf diese Weise werden die Folgen $m,m',m'',m'''$ usw. und $n,n',n'',n'''$ usw. sehr rasch gegen einen \emph{gemeinsamen Grenzwert} konvergieren, den wir mit $\mu$ bezeichnen und einfach das \emph{arithmetisch-geometrische Mittel} zwischen $m$ und $n$ nennen werden. Wir werden nun beweisen, dass $\frac1\mu$ der Wert des Integrals
\[
\begin{aligned}
&\int\frac{\mathrm dT}{2\pi\surd{(mm\cos T^2+nn\sin T^2)}}
\end{aligned}
\]
ist, erstreckt von $T=0$ bis $T=360^0$.

\Needspace{40\baselineskip}
\emph{Beweis.} Wir nehmen an, dass die Variable $T$ durch eine andere $T'$ so ausgedrückt werde, dass
\[
\begin{aligned}
&\sin T=\frac{2m\sin T'}{(m+n)\cos T^{\prime 2}+2m\sin T^{\prime 2}}
\end{aligned}
\]
und man wird leicht erkennen, dass, während $T'$ vom Wert $0$ bis zu $90^0,180^0,270^0,360^0$ zunimmt, auch $T$, wenn auch in ungleichen Intervallen, von $0$ bis zu $90^0,180^0,270^0,360^0$ wächst. Nach gehöriger Entwicklung findet man
\[
\begin{aligned}
&\frac{\mathrm dT}{\surd{(mm\cos T^2+nn\sin T^2)}}
=\frac{\mathrm dT'}{\surd{(m'm'\cos T^{\prime 2}+n'n'\sin T^{\prime 2})}}
\end{aligned}
\]
\SourceLeaf{362}
\Needspace{34\baselineskip}
und somit sind die Werte der Integrale
\[
\begin{gathered}\int\frac{\mathrm{d}T}{2\pi\surd(mm\cos T^2+nn\sin T^2)},\\[7pt]
\int\frac{\mathrm{d}T'}{2\pi\surd(m'm'\cos T'^2+n'n'\sin T'^2)}\end{gathered}
\]
einander gleich, wenn jede der beiden Variablen vom Wert $0$ bis zum Wert $360^{0}$ erstreckt wird. Da man auf dieselbe Weise weiter fortfahren darf, ist klar, dass diesen Werten auch der Wert des Integrals
\[
\int\frac{\mathrm{d}\Theta}{2\pi\surd(\mu\mu\cos\Theta^2+\mu\mu\sin\Theta^2)}
\]
von $\Theta=0$ bis $\Theta=360^{0}$ gleich ist; dieser wird offenbar $=\frac{1}{\mu}$. Q.~E.~D.

\Needspace{44\baselineskip}
\AttractionArticle{17}
Aus der Gleichung, welche die Beziehung zwischen $T$ und $T'$ angibt,
\[
(m-n)\sin T.\sin T'^2=2m\sin T'-(m+n)\sin T
\]
folgt leicht
\[
\begin{aligned}
\surd(mm\cos T^2+nn\sin T^2)&=m-(m-n)\sin T.\sin T'\\[6pt]
\surd(m'm'\cos T'^2+n'n'\sin T'^2)&=m\mathop{\mathrm{cotang}}T.\mathop{\mathrm{tang}}T'
\end{aligned}
\]
und hieraus mit Hilfe derselben Gleichung
\[
\begin{aligned}
&\sin T.\sin T'.\surd(mm\cos T^2+nn\sin T^2)+m'(\cos T^2-\sin T^2)\\[6pt]
&=\cos T.\cos T'.\surd(m'm'\cos T'^2+n'n'\sin T'^2)-\tfrac12(m-n)\sin T'^2
\end{aligned}
\]
\par Multipliziert man diese Gleichung mit
\[
\frac{\mathrm{d}T}{\surd(mm\cos T^2+nn\sin T^2)}
=\frac{\mathrm{d}T'}{\surd(m'm'\cos T'^2+n'n'\sin T'^2)}
\]
so ergibt sich
\[
\begin{aligned}&\frac{m'(\cos T^2-\sin T^2)\mathrm{d}T}{\surd(mm\cos T^2+nn\sin T^2)}\\[7pt]&=-\frac{\tfrac12(m-n)\sin T'^2\mathrm{d}T'}{\surd(m'm'\cos T'^2+n'n'\sin T'^2)}+\mathrm{d}.\sin T'\cos T\end{aligned}
\]
\Needspace{30\baselineskip}
\par Multiplizieren wir diese Gleichung mit $\frac{m-n}{\pi}$, setzen wir
$m'(m-n)=\tfrac12(mm-nn)$,
$(m-n)^2=4(m'm'-n'n')$,
$\sin T'^2=\tfrac12-\tfrac12(\cos T'^2-\sin T'^2)$ ein und integrieren von den Werten $T$ und $T'=0$ bis $360^{0}$, so haben wir:
\[
\begin{aligned}
&(mm-nn)\int\frac{(\cos T^2-\sin T^2).\mathrm{d}T}{2\pi\surd(mm\cos T^2+nn\sin T^2)}\\[6pt]
&=-\frac{2(m'm'-n'n')}{\mu}\\[7pt]&\quad+2(m'm'-n'n')\int\frac{(\cos T'^2-\sin T'^2)\mathrm{d}T'}{2\pi\surd(m'm'\cos T'^2+n'n'\sin T'^2)}
\end{aligned}
\]
\SourceLeaf{363}
Da sich das bestimmte Integral auf der rechten Seite auf dieselbe Weise umformen lässt, wird das Integral
\[
\int\frac{(\cos T^2-\sin T^2)\mathrm{d}T}{2\pi\surd(mm\cos T^2+nn\sin T^2)}
\]
offenbar durch die sehr rasch konvergierende unendliche Reihe
\[
-\frac{\begin{gathered}2(m'm'-n'n')+4(m''m''-n''n'')\\+8(m'''m'''-n'''n''')+\mathrm{usw.}\end{gathered}}{(mm-nn)\mu}
=-\frac{\nu}{\mu}
\]
ausgedrückt werden.
\Needspace{31\baselineskip}
\par Die numerische Rechnung wird am bequemsten durch Logarithmen ausgeführt, wenn wir
\[
\begin{gathered}\tfrac14\surd(mm-nn)=\lambda,\quad\tfrac14\surd(m'm'-n'n')=\lambda',\\[6pt]\tfrac14\surd(m''m''-n''n'')=\lambda''\ \mathrm{usw.}\end{gathered}
\]
setzen; daraus wird sich ergeben:
\[
\lambda'=\frac{\lambda\lambda}{m'},\quad
\lambda''=\frac{\lambda'\lambda'}{m''},\quad
\lambda'''=\frac{\lambda''\lambda''}{m'''}\ \mathrm{usw.}\quad\text{und}
\]
\[
\nu=\frac{2\lambda'\lambda'+4\lambda''\lambda''+8\lambda'''\lambda'''+\mathrm{usw.}}{\lambda\lambda}
\]

\clearpage
\endgroup
